\documentclass[11pt]{amsart}
\usepackage{geometry}                % See geometry.pdf to learn the layout options. There are lots.
\geometry{letterpaper}                   % ... or a4paper or a5paper or ... 
%\geometry{landscape}                % Activate for for rotated page geometry
%\usepackage[parfill]{parskip}    % Activate to begin paragraphs with an empty line rather than an indent
\usepackage{graphicx}
\usepackage{amssymb}
\usepackage{epstopdf}
\DeclareGraphicsRule{.tif}{png}{.png}{`convert #1 `dirname #1`/`basename #1 .tif`.png}

\title{PS 4.22}
%\date{}                                           % Activate to display a given date or no date

\begin{document}
\maketitle
%\section{}
%\subsection{}
 \noindent Los datos reunidos en la Secretaría de Industrias relacionado a la producción de N fábricas en cada uno de los meses del año anterior, se proporcionan de la siguiente manera:  \newline
Datos: N  \newline
\begin{center}
$FAB_1$, $MES_1,_1$, $MES_1,_2$, $MES_1,_{12}$,  \newline
$FAB_2$, $MES_2,_1$, $MES2,_2$, $MES_2,_{12}$,  \newline
...  \newline
... \newline
 \noindent$FAB_N$, $MES_N,_1$, $MES_N,_2$, $MES_N,_{12}$
\end{center} 

Donde:  

\setlength{\parindent}{1cm}
N \hspace{0.1cm} es una variable de tipo entero que representa el número de fábricas,      $1 \leq N \leq 500$ 

\setlength{\parindent}{1cm}
$FAB_i$ \hspace{0.5cm} es una variable de tipo entero que representa la clave de la fábrica$_i$ 

\setlength{\parindent}{1cm}
$MES_i,_j$  \hspace{0.5cm} es una variable de tipo real que representa la producción de la fábrica$_j$ en el mes j.  \newline

 \noindent Realice un diagrama de flujo que proporcione la siguiente información:
\begin{enumerate}
  \item[a)] La clave de la fábrica que más produjo el año anterior. Imprimir también su producción.
  \item[b)] Dado un mes como parámetro, imprimir las claves de las fábricas cuyas producciones en dicho mes, fueron superiores a \$150,000.
\end{enumerate}
\end{document}  